%STEPBEGIN{0}
\documentclass{standalone}
\usepackage{hsbcommon}
\begin{document}
\showmytreeraw[5]{forest-loft-bed.tex}   %%% <-- EDIT step list; keep this filename in sync with the file's own name
\end{document}
%STEPEND
% Step 1: root only
% Step 2: + Execution / Preparation
% Step 3: + Preparation's tasks
% Step 4: + Execution's tasks
% Step 5: + "make a plan" broken down
% Where a node that later becomes a container has to be hidden, its opening
% AND closing bracket are tagged with the SAME step spec, so every step's
% filtered text stays balanced. Tags are never nested.
    \begin{forest}
    for tree={
        font=\footnotesize,
        grow = east,
        anchor = west,
        parent anchor = east,
        child anchor = west,
        forked edge,
        inner sep=0.5pt
            }
    [Build a loft bed
        %STEPBEGIN{2,3,4,5}
        [Execution
        %STEPEND
            %STEPBEGIN{4,5}
            [assemble]
            [transport]
            [paint or conserve]
            [shape materials]
            [rent space and tools]
            [order material]
            %STEPEND
        %STEPBEGIN{2,3,4,5}
        ]
        [Preparation
        %STEPEND
            %STEPBEGIN{3,4,5}
            [find helpers]
            [order material]
            [organise tools and workshop]
            [finance]
            [make a plan
            %STEPEND
                %STEPBEGIN{5}
                [calculate]
                [draw]
                [measure]
                [consult]
                %STEPEND
            %STEPBEGIN{3,4,5}
            ]
            %STEPEND
        %STEPBEGIN{2,3,4,5}
        ]
        %STEPEND
    ]
    \end{forest}